% =============================================================================
%  DeliveryBench-RL --- paper outline / starting point
%
%  Self-contained `article` skeleton so it compiles anywhere:
%      pdflatex deliverybench_rl.tex   (run twice for refs)
%  To submit to a venue, swap the documentclass/preamble for the conference
%  style (e.g. \documentclass{article}\usepackage{neurips_2025}) and keep the
%  body below.
%
%  Prose is written out where the content is settled (Introduction,
%  DeliveryBench 2.0). Experiments are placeholders with fully specified
%  figures/tables --- search for \TODO and the "PLACEHOLDER" boxes.
% =============================================================================
\documentclass[11pt]{article}

\usepackage[utf8]{inputenc}
\usepackage[T1]{fontenc}
\usepackage[margin=1in]{geometry}
\usepackage{amsmath,amssymb}
\usepackage{graphicx}
\usepackage{booktabs}
\usepackage{multirow}
\usepackage{array}
\usepackage{xcolor}
\usepackage{pifont}      % \ding for check / cross marks
\usepackage{enumitem}
\usepackage{subcaption}
\usepackage{microtype}
\usepackage[colorlinks=true,allcolors=blue!55!black]{hyperref}
\usepackage[capitalize,noabbrev]{cleveref}

% ---- shorthands -------------------------------------------------------------
\newcommand{\ours}{\textsc{DeliveryBench-RL}\xspace}
\usepackage{xspace}
\newcommand{\dbone}{\textsc{DeliveryBench}\xspace}
\newcommand{\TODO}[1]{{\color{red}\textbf{[TODO:} #1\textbf{]}}}
\newcommand{\cmark}{\textcolor{green!55!black}{\ding{51}}}      % full support
\newcommand{\xmark}{\textcolor{red!70!black}{\ding{55}}}        % no support
\newcommand{\pmark}{\textcolor{orange!85!black}{$\sim$}}        % partial support
\newcommand{\unit}[1]{\textsf{#1}}                              % a unit capability

% gray placeholder box for figures not yet rendered
\newcommand{\figplaceholder}[2][0.85\linewidth]{%
  \begin{center}
  \fbox{\begin{minipage}[c][#2][c]{#1}\centering
    \color{gray}\itshape Figure placeholder --- render from results.
  \end{minipage}}
  \end{center}}

\title{\ours: Reinforcement Learning for Long-Horizon,\\
       Multimodal Embodied Agents under Realistic Constraints}
\author{\TODO{author list} \\ \TODO{affiliation}}
\date{}

\begin{document}
\maketitle

% =============================================================================
\begin{abstract}
Coding agents have shown that strong autonomy, tool use, and proactive
reasoning can be elicited through long-horizon harnesses and agentic training.
We ask whether the same agentic capacity can be granted to a \emph{real-world,
embodied} setting with multimodal input. \dbone introduced a long-horizon
embodied delivery environment with realistic constraints---time, energy, and
safety---and revealed two core deficiencies that persist even in frontier
vision--language models (VLMs): a lack of long-horizon planning, and a failure
to \emph{proactively} satisfy task constraints, including resource limits and
\emph{implicit} commonsense rules (e.g.\ safety and traffic rules) that the
same models recall reliably in isolated question answering.
We present \ours, which turns this benchmark into a trainable environment for
multimodal agentic reinforcement learning (RL) by adding step-wise reward
shaping, an oracle planner, a multi-dimensional curriculum, tool calling, and
adjustable action granularity. We train and compare RL baselines
(PPO, GRPO, SUPO) against prompt-scaffolded frontier VLMs, and dissect the
extent to which RL grants proactive mastery of individual implicit constraints.
\TODO{1--2 sentence headline result.}
\end{abstract}

% =============================================================================
\section{Introduction}
\label{sec:intro}

Recent progress in \emph{coding agents} has shown that strong autonomy, fluent
tool use, and proactive, self-directed reasoning can be \emph{elicited and
trained} rather than merely prompted, given a long-horizon harness and an
agentic training loop. A natural question follows: can the same agentic
capacity be granted in a \emph{real-world embodied} environment with
\emph{multimodal} input, where an agent must perceive, plan over long horizons,
and respect physical and social constraints?

\dbone was, to our knowledge, the first long-horizon agentic environment for
real-world embodied tasks to couple this multimodal, egocentric setting with
\emph{realistic constraints}---a time budget (a fixed shift), energy and
battery limits, and safety/social rules such as obstacles and traffic signals.
Studying frontier VLMs in \dbone exposed two core deficiencies that survive
even at the frontier:
\begin{enumerate}[leftmargin=1.4em,itemsep=2pt,topsep=2pt]
  \item \textbf{Weak long-horizon planning.} Models fail to compose the many
        steps of a delivery (select, route, pick up, deliver, replan) into a
        plan that respects a global time/energy budget.
  \item \textbf{Non-proactive constraint handling.} Models do not
        \emph{proactively} use task constraints---both resource limits and
        \emph{implicit} commonsense rules---to guide their reasoning. Crucially,
        these rules (e.g.\ ``do not cross on red'', ``go around an obstacle'')
        are recalled reliably in \emph{direct question answering}, yet when the
        same rule must be applied from within a complex, long-horizon task
        state, the model fails to invoke it. This gap between
        \emph{knowing} a rule and \emph{acting} on it is a latent robustness and
        safety problem.
\end{enumerate}

\ours builds directly on this foundation and makes \emph{multimodal agentic RL
training} in the environment feasible. We integrate (i) step-wise
\textbf{reward shaping} with a customizable, environment-aware interface;
(ii) an \textbf{oracle planner} that yields globally optimal trajectories;
(iii) a multi-dimensional \textbf{curriculum}; (iv) \textbf{tool calling} for
adaptive context loading; and (v) adjustable \textbf{action granularity}.
The oracle planner in particular gives us infrastructure for imitation
learning, reward normalization, automatic QA/unit-test construction, and
fine-grained failure analysis, while action granularity lets the same
environment probe high-level planning \emph{and} low-level visual navigation
(spatial reasoning, safety, social navigation).

\paragraph{Contributions.}
\begin{enumerate}[leftmargin=1.4em,itemsep=2pt,topsep=2pt]
  \item We extend \dbone into \ours, an RL-ready, long-horizon, multimodal
        embodied environment with an oracle planner, multi-dimensional
        curriculum, step-wise (and customizable) reward shaping, tool calling,
        and adjustable action granularity (\cref{sec:env}).
  \item We conduct, to our knowledge, the first study of multimodal agentic RL
        in a real-world long-horizon embodied environment, training and
        comparing baseline algorithms (PPO, GRPO, SUPO) and prompt-scaffolded
        frontier VLMs (\cref{sec:exp-main}).
  \item We analyze how well RL-trained models acquire \emph{proactive} handling
        of distinct implicit constraints, and isolate the \emph{unit}
        capabilities models still struggle to master (\cref{sec:exp-units}).
\end{enumerate}

% =============================================================================
\section{Related Work}
\label{sec:related}
\TODO{Long-horizon / agentic RL environments; embodied \& navigation
benchmarks (ALFWorld, Habitat, etc.); LLM/VLM agents and tool use; agentic RL
(PPO/GRPO and variants); curriculum learning; imitation from oracle/expert.
Position \ours against \cref{tab:env-comparison}.}

% =============================================================================
\section{\ours: An RL-Ready Embodied Agentic Environment}
\label{sec:env}

\ours starts from \dbone's simulated city-delivery task---a courier completes
food deliveries within a two-hour shift under time, energy/battery, and
safety constraints---and adds the machinery required to \emph{train} agents,
not just evaluate them. \Cref{tab:env-comparison} contrasts \ours with \dbone
and with other long-horizon agentic environments. The five additions below are
ordered by their centrality to enabling RL.

% ---- comparison table -------------------------------------------------------
\begin{table}[t]
\centering
\caption{\ours versus \dbone and other long-horizon agentic environments.
\cmark~full support, \pmark~partial, \xmark~none. The top block is the direct
\dbone lineage; the bottom block (other environments) is \TODO{populate with
verified entries}.}
\label{tab:env-comparison}
\setlength{\tabcolsep}{4pt}
\renewcommand{\arraystretch}{1.15}
\resizebox{\textwidth}{!}{%
\begin{tabular}{l ccc ccccc c}
\toprule
 & \rotatebox{55}{Multimodal (FPV+map)}
 & \rotatebox{55}{Long-horizon}
 & \rotatebox{55}{Realistic constraints}
 & \rotatebox{55}{Oracle planner}
 & \rotatebox{55}{Curriculum}
 & \rotatebox{55}{Step-wise reward shaping}
 & \rotatebox{55}{Tool calling}
 & \rotatebox{55}{Adjustable action granularity}
 & \rotatebox{55}{RL-ready} \\
\midrule
\textbf{\ours (ours)} & \cmark & \cmark & \cmark & \cmark & \cmark & \cmark & \cmark & \cmark & \cmark \\
\dbone (v1)           & \cmark & \cmark & \cmark & \xmark & \xmark & \pmark & \xmark & \pmark & \xmark \\
\midrule
\multicolumn{10}{l}{\emph{Other long-horizon agentic environments}}\\
\TODO{Env A} & \TODO{} & \TODO{} & \TODO{} & \TODO{} & \TODO{} & \TODO{} & \TODO{} & \TODO{} & \TODO{} \\
\TODO{Env B} & \TODO{} & \TODO{} & \TODO{} & \TODO{} & \TODO{} & \TODO{} & \TODO{} & \TODO{} & \TODO{} \\
\TODO{Env C} & \TODO{} & \TODO{} & \TODO{} & \TODO{} & \TODO{} & \TODO{} & \TODO{} & \TODO{} & \TODO{} \\
\bottomrule
\end{tabular}}
\end{table}

% ---- 3.1 reward shaping -----------------------------------------------------
\subsection{Reward Shaping: From Terminal Profit to Expressive Step-wise Signal}
\label{sec:env-reward}
In \dbone the only reward signal was the \emph{terminal} profit of a full
trajectory---a sparse, high-variance target that is poorly suited to
long-horizon RL. \ours decomposes this into a \textbf{step-wise} reward built
from short-horizon \emph{partial-success} events: response/format-schema
correctness, constraint violations (e.g.\ collisions, red-light crossings,
deadline misses), and sub-task completions (order accepted, picked up, dropped
off). This makes the training signal far more expressive and densely
attributable. Importantly, we expose an \textbf{interface for customizing
reward shaping from environment state}, so downstream work can define new
shaping terms (or safety penalties) without touching the core loop.
\TODO{State the default reward decomposition $r_t = \sum_k w_k\, \phi_k(s_t,a_t)$
and list the terms $\phi_k$ in a table; note which are on by default.}

% ---- 3.2 oracle planner -----------------------------------------------------
\subsection{Oracle Planner: Optimal Trajectories for Normalization, Imitation, and Tests}
\label{sec:env-oracle}
Using full global information, a search-based planner computes a
(near-)optimal trajectory for any task instance. This single component powers
four capabilities: (i) \textbf{reward normalization}---scoring an agent against
the optimal achievable profit/cost rather than an absolute scale;
(ii) \textbf{imitation learning}---oracle rollouts become expert demonstrations
for behavior cloning / SFT warm-start; (iii) \textbf{automatic QA and unit
tests}---each oracle decision yields a labeled ``what is the correct next
action / which constraint applies here'' probe; and (iv) \textbf{failure
analysis}---an agent trajectory can be diffed against the oracle to localize
\emph{where} and \emph{why} it diverged.
\TODO{Specify the planner (graph search over the waypoint graph; objective =
maximize profit / minimize time--energy s.t.\ deadlines and safety); note
optimality guarantees / approximations.}

% ---- 3.3 curriculum ---------------------------------------------------------
\subsection{Multi-Dimensional Curriculum}
\label{sec:env-curriculum}
\ours warms the agent into the environment along three orthogonal axes:
\begin{itemize}[leftmargin=1.4em,itemsep=2pt,topsep=2pt]
  \item \textbf{Horizon curriculum}: progressively longer shifts and more
        orders per episode, growing the planning horizon.
  \item \textbf{Prompt curriculum}: the description of each constraint and the
        action space is annealed from \emph{few-shot} examples, to explicit
        \emph{instructions}, to \emph{proactive} (the agent must surface the
        rule itself)---directly targeting the know-vs-act gap of
        \cref{sec:intro}.
  \item \textbf{Constraint curriculum}: realistic constraints and their
        associated actions are introduced incrementally (e.g.\ energy/battery,
        then obstacles/\textsc{bypass}, then traffic lights/\textsc{wait}).
\end{itemize}
\TODO{Table of curriculum stages $\times$ axis settings; which stage unlocks
which constraint/action.}

% ---- 3.4 tool calling -------------------------------------------------------
\subsection{Tool Calling for Adaptive Context Loading}
\label{sec:env-tools}
Rather than dumping every observation channel into the step prompt, \ours
exposes the four-direction first-person views, the top-down map, and
\emph{query tools} (navigation/route, order pool, energy/battery, \dots) that
the agent calls on demand. This \textbf{adaptive context loading} keeps each
step's context small, which is what makes long-horizon rollouts tractable to
\emph{scale}---context grows with what the agent chooses to inspect, not
linearly with the horizon.
\TODO{List the tools and their costs (if any); note which are query-only vs
state-changing.}

% ---- 3.5 action granularity -------------------------------------------------
\subsection{Adjustable Action Granularity}
\label{sec:env-granularity}
Finally, \ours makes action granularity a configurable axis. At a coarse
granularity the environment probes high-level \emph{long-horizon planning and
constraint understanding}; at a fine granularity, navigation is driven directly
from the egocentric FPV and map, letting the same environment evaluate
\emph{spatial reasoning}, \emph{safety}, and \emph{social navigation}
(e.g.\ noticing a road obstacle in the front view and choosing to go around it,
or waiting at a red light). This is what lets a single benchmark span planning
and perception-grounded control.
\TODO{Define the granularity levels (e.g.\ macro \textsc{navigate} vs primitive
\textsc{move(direction)}) and what each evaluates.}

% =============================================================================
\section{Experiments}
\label{sec:exp}

\paragraph{Setup.}
\TODO{Backbone open-source VLM (e.g.\ Qwen3-VL-8B), frontier VLMs for the
scaffolded baselines, RL framework, train/eval splits, seeds, compute,
evaluation metrics (success / deliveries, hourly profit, constraint-violation
rates, deadline misses), and which curriculum schedule is used.}

% ---- 4.1 main table ---------------------------------------------------------
\subsection{Main Results: Frontier Scaffolding vs.\ RL Baselines}
\label{sec:exp-main}

\begin{table}[t]
\centering
\caption{Main results. Prompt-scaffolded frontier VLMs vs.\ open-source RL
baselines on \ours. Metrics: task success / \#deliveries, hourly profit
(profit $/$ sim-hour), and constraint-violation rates (collisions, traffic
violations, deadline misses); lower is better for violations.
\TODO{fill all cells; bold the best per column; add std over seeds}.}
\label{tab:main}
\setlength{\tabcolsep}{6pt}
\renewcommand{\arraystretch}{1.1}
\resizebox{\textwidth}{!}{%
\begin{tabular}{l l ccc ccc}
\toprule
 & & \multicolumn{3}{c}{Task performance} & \multicolumn{3}{c}{Constraint violations ($\downarrow$)}\\
\cmidrule(lr){3-5}\cmidrule(lr){6-8}
Method & Type & Success & \#Deliv. & \$/h & Collisions & Traffic & Deadline \\
\midrule
\multicolumn{8}{l}{\emph{Prompt-scaffolded frontier VLMs (no training)}}\\
\TODO{GPT-x}      & zero/few-shot & \TODO{} & \TODO{} & \TODO{} & \TODO{} & \TODO{} & \TODO{} \\
\TODO{Claude}     & zero/few-shot & \TODO{} & \TODO{} & \TODO{} & \TODO{} & \TODO{} & \TODO{} \\
\TODO{Gemini}     & zero/few-shot & \TODO{} & \TODO{} & \TODO{} & \TODO{} & \TODO{} & \TODO{} \\
\midrule
\multicolumn{8}{l}{\emph{Open-source backbone + RL (ours)}}\\
\TODO{Qwen3-VL-8B} & base (no RL)  & \TODO{} & \TODO{} & \TODO{} & \TODO{} & \TODO{} & \TODO{} \\
\quad + PPO        & RL            & \TODO{} & \TODO{} & \TODO{} & \TODO{} & \TODO{} & \TODO{} \\
\quad + GRPO       & RL            & \TODO{} & \TODO{} & \TODO{} & \TODO{} & \TODO{} & \TODO{} \\
\quad + SUPO       & RL            & \TODO{} & \TODO{} & \TODO{} & \TODO{} & \TODO{} & \TODO{} \\
\bottomrule
\end{tabular}}
\end{table}

\Cref{tab:main} compares prompt-scaffolded frontier VLMs against an open-source
backbone trained with PPO, GRPO, and SUPO. \TODO{1--3 sentences of takeaway:
does RL on a small open model close / exceed the gap to scaffolded frontier
models, and on which metric (planning/profit vs.\ safety/violations)?}

% ---- 4.2 curriculum emergence ----------------------------------------------
\subsection{Curriculum Dynamics: Staged Emergence of Unit Tasks}
\label{sec:exp-curriculum}

\begin{figure}[t]
\figplaceholder{6cm}
\caption{\textbf{Reward / per-unit-task success vs.\ training steps}, annotated
with curriculum-stage boundaries (vertical dashed lines).
\emph{What to plot:} x-axis = environment / training steps; left y-axis =
mean episode reward (or hourly profit); overlaid curves = success rate of each
\emph{unit task} (\unit{accept}, \unit{navigate}, \unit{pickup},
\unit{dropoff}, \unit{energy-management}, \unit{obstacle-avoidance},
\unit{traffic-compliance}). The figure should show each unit capability
\emph{emerging} in turn as the corresponding curriculum stage unlocks, rather
than all at once. Optionally add an ablation curve (no curriculum) to show
slower / non-emergent learning.}
\label{fig:curriculum}
\end{figure}

\Cref{fig:curriculum} tracks reward and per-unit-task success across training.
\TODO{Takeaway: which unit tasks emerge early vs.\ late; whether the curriculum
is necessary for the hardest units (e.g.\ proactive safety) to emerge at all.}

% ---- 4.3 unit tests / proactivity ------------------------------------------
\subsection{Unit Tests: Proactive Mastery of Implicit Constraints}
\label{sec:exp-units}

\begin{figure}[t]
\figplaceholder{6cm}
\caption{\textbf{Direct-QA accuracy vs.\ in-task proactive use, per constraint /
subtask.} \emph{What to plot:} grouped bars (or a radar) over constraints and
subtasks---\unit{obstacle/safety}, \unit{traffic-light}, \unit{energy budget},
\unit{deadline}, \unit{packing}, \unit{mode selection}, \dots---with two bars
per group: (a) accuracy when the rule is asked in \emph{isolated QA}, and
(b) the rate at which the model \emph{proactively} applies the rule from within
a live long-horizon state (extracted via the oracle unit tests of
\cref{sec:env-oracle}). The \emph{gap} (a)$-$(b) quantifies the
know-vs-act deficiency; show it for the base model and after RL.}
\label{fig:units}
\end{figure}

\Cref{fig:units} isolates, per constraint, the gap between \emph{knowing} a rule
(direct QA) and \emph{proactively acting} on it inside a long-horizon task.
\TODO{Takeaway: which implicit constraints RL teaches the model to use
proactively, and which it still fails on (the residual unit capabilities);
connect back to the safety/robustness framing in \cref{sec:intro}. Reference
the obstacle-avoidance probe (e.g.\ proactive \textsc{bypass} vs.\ collision)
as a concrete instance.}

% =============================================================================
\section{Analysis and Discussion}
\label{sec:analysis}
\TODO{Failure-mode taxonomy via oracle diffing (\cref{sec:env-oracle}); where RL
helps (planning/profit) vs.\ where it lags (proactive safety); effect of action
granularity on spatial reasoning; limitations (sim-to-real, single map,
reward-hacking risks); broader impact (safety of embodied agents).}

% =============================================================================
\section{Conclusion}
\label{sec:conclusion}
\TODO{Restate: \ours makes long-horizon multimodal embodied RL feasible; RL
baselines close part of the gap to scaffolded frontier VLMs but proactive
implicit-constraint handling remains an open unit capability.}

% =============================================================================
\appendix
\section{Environment Details}
\TODO{Action/observation spaces, constraint definitions and costs, reward-term
table, curriculum schedule, tool specifications, oracle-planner algorithm.}

\section{Training Details}
\TODO{Hyperparameters per algorithm (PPO/GRPO/SUPO), context management,
compute, wall-clock, seeds.}

% \bibliographystyle{plainnat}
% \bibliography{references}

\end{document}
